\section{Significance of the human-fed blind spot across independent interventions}
\label{sec:significance}

\subsection{A single-condition test}

Benchmark~2 evaluated agency accuracy on the $n=17$ clips with true label
\texttt{human-fed} under several independent elicitation conditions:

\begin{table}[h]
\centering
\begin{tabular}{lr}
\toprule
Condition & Agency accuracy, true human-fed ($n=17$) \\
\midrule
Composite verb question                & $0\%$ \\
Agency asked alone                      & $0\%$ \\
``Whose hand holds the food?'' (perception probe) & $6\%$ \\
Describe-people-first scaffold          & $18\%$ \\
``fed-by-another-person'' rewording     & $0\%$ \\
\bottomrule
\end{tabular}
\caption{Reproduced from Benchmark 2 / Benchmark 5.}
\label{tab:agency-conditions}
\end{table}

Fix a null hypothesis $H_0$: the model's per-trial probability of correctly
identifying agency on a true human-fed clip is some baseline rate $p_0$
representative of ``the model can perceive agency at least as well as it
does on the easier self-fed subgroup'' (empirically $\approx 90\%$ there),
or, more conservatively, $p_0 = 0.5$ (``no better than an agency coin
flip''). Under $H_0$, with $K$ the number of correct calls out of $n=17$
i.i.d.\ trials, $K \sim \mathrm{Binomial}(17, p_0)$, and the observed
$K=0$ has one-sided p-value
\[
  \Pr(K \le 0 \mid n=17, p_0) \;=\; (1-p_0)^{17}.
\]
\[
  p_0 = 0.5: \quad (0.5)^{17} \approx 7.6\times 10^{-6}, \qquad
  p_0 = 0.9: \quad (0.1)^{17} \approx 10^{-17}.
\]
Either way, observing zero correct calls out of 17 is not plausibly
explained by ordinary sampling noise around a reasonable baseline rate ---
it is evidence the true per-trial success probability on this subgroup is
near zero, not merely below its self-fed-subgroup counterpart.

\subsection{Combining across independent conditions: Fisher's method}

Table~\ref{tab:agency-conditions} contains five \emph{semi-independent}
probes (independent in the sense that each perturbs a different aspect of
elicitation --- question framing, isolated sub-question, pure perception,
reasoning scaffold, label vocabulary --- rather than being repeated trials
of the identical prompt). Let $p_i$ be the one-sided binomial p-value
computed as above (under $H_0{:}\,p_0=0.5$) for condition $i \in
\{1,\dots,5\}$, using the reported accuracy to back out an equivalent
integer success count $K_i = \mathrm{round}(0.17 \times \mathrm{acc}_i)$
out of $n=17$:\footnote{$K_i = 0$ for the three $0\%$ conditions, $K_i=1$
for the $6\%$ probe, $K_i=3$ for the $18\%$ scaffold.}

\begin{definition}[Fisher's combined probability test]
\[
  \chi^2_{\mathrm{combined}} \;=\; -2\sum_{i=1}^{5} \ln p_i \;\sim\; \chi^2_{2\times 5 = 10} \quad \text{under } H_0.
\]
\end{definition}

With $p_1=p_2=p_5 = \Pr(K\le 0\mid n{=}17,p_0{=}0.5) \approx 7.63\times10^{-6}$
(the three $0\%$ conditions),
$p_3 = \Pr(K\le 1\mid n{=}17,p_0{=}0.5) \approx 1.37\times10^{-4}$, and
$p_4 = \Pr(K\le 3\mid n{=}17,p_0{=}0.5) \approx 6.36\times10^{-3}$,
\[
  \chi^2_{\mathrm{combined}} \;\approx\; -2\big(3\ln(7.63\times10^{-6}) + \ln(1.37\times10^{-4}) + \ln(6.36\times10^{-3})\big) \;\approx\; 98.6,
\]
against a $\chi^2_{10}$ null distribution whose $99.9$th percentile is
about $29.6$ --- the combined statistic exceeds it by more than $3\times$,
giving a combined p-value (numerically, $\mathrm{sf}_{\chi^2_{10}}(98.6)
\approx 1.0\times10^{-16}$) many orders of magnitude below any conventional
significance threshold.

\begin{remark}[What this rules in and rules out]
This computation does not claim the five conditions are statistically
independent in the strict i.i.d.\ sense (they share the same 17 clips and
the same underlying model weights), so the nominal $\chi^2_{10}$
$p$-value should be read qualitatively rather than taken at face value.
What it does formalize is the qualitative argument used in
Section~\ref{sec:reliability}: a fixable \emph{prompting} deficiency
predicts that \emph{at least one} of five substantially different framings
(especially the pure-perception probe, which asks nothing about the
taxonomy at all) should recover non-trivial accuracy, whereas a
\emph{perceptual} deficiency in how the visual encoder or attention
mechanism attributes a hand to a person predicts exactly the observed
pattern: uniformly near-zero accuracy regardless of how the question is
phrased, with only the scaffold condition --- which adds an explicit
step of \emph{describing} what is seen before answering --- producing a
partial (and still far from adequate) recovery to $18\%$.
\end{remark}

\subsection{Interaction with the few-shot result}

Benchmark~5, Part B adds a sixth and seventh condition (4-shot and 8-shot
in-context demonstrations, the latter explicitly including two human-fed
examples) and observes $0\%$ agency accuracy on true human-fed clips at
\emph{both} shot counts. Extending the binomial argument, providing
correctly labeled human-fed exemplars in-context is the single
intervention most directly targeted at a demonstration-availability
explanation of the failure (``the model has never been shown what
human-fed looks like in this schema''); its failure to move $K$ away from
$0$ is the strongest single piece of evidence in the suite that the
deficit is not remediable by prompting or in-context learning at the
current model scale, and would require either fine-tuning on labeled
human-fed examples or an architectural/attention-level intervention (e.g.\
person-conditioned attention, as sketched by the scaffold condition's
partial success) to resolve.
